\subsection{Case III: CWRU Bearing Dataset}
The Case III evaluation is conducted on the Case Western Reserve University (CWRU) bearing dataset \cite{Smith2015CWRU}, which serves as a benchmark for assessing the robustness and generalization of fault diagnosis algorithms under varying operating conditions. The dataset consists of vibration signals collected from the drive end of a motor-driven system at sampling rates of 12~kHz and 48~kHz. The experimental setup covers ten distinct health conditions, including one normal state and nine fault states. Specifically, the faults are categorized into inner raceway (IR), outer raceway (OR), and ball (BA) faults, each manifesting in three different damage diameters (0.007, 0.014, and 0.021 inch). Furthermore, data were recorded under four motor loads ranging from 0 to 3 hp, introducing significant domain shifts that challenge the zero-shot learning (ZSL) capabilities of the proposed Diff-LM-GZSL model.

To implement the zero-shot fault diagnosis task, the dataset is partitioned into seven seen fault types used during training and three unseen fault types reserved exclusively for testing. To leverage the powerful image-generation capabilities of the latent diffusion module, the raw one-dimensional time-series vibration signals are transformed into two-dimensional time-frequency representations. This is achieved using the Short-Time Fourier Transform (STFT) and Continuous Wavelet Transform (CWT), producing 2D spectrograms that encapsulate both temporal and spectral characteristics of the bearing dynamics. These spectrograms provide a rich visual structure for the diffusion model to learn the underlying data distribution conditioned on the linguistic semantic embeddings generated by the LLM.

Table \ref{tab:cwru_results} presents the performance comparison between Diff-LM-GZSL and nine state-of-the-art baselines on the CWRU dataset. The metrics include the seen class accuracy ($acc\_s$), unseen class accuracy ($acc\_u$), harmonic mean ($H$), and F1-score. As shown in the table, Diff-LM-GZSL achieves a harmonic mean of 89.2\%, significantly outperforming the second-best baseline, which scores 82.4\%. The high $acc\_u$ (87.5\%) indicates that the proposed method effectively transfers knowledge from the seen classes to the unseen categories by aligning the latent diffusion space with the high-dimensional linguistic semantics. The standard deviations across five independent runs are maintained within $\pm$0.6\%, demonstrating the stability of the diffusion-based generation process compared to GAN-based alternatives.

\begin{table}[ht]
\centering
\caption{Performance Comparison on CWRU Bearing Dataset (Case III)}
\label{tab:cwru_results}
\begin{tabular}{lcccc}
\hline
Method & $acc\_s$ (\%) & $acc\_u$ (\%) & $H$ (\%) & F1-score \\ \hline
CNN-Softmax & 98.2 $\pm$ 0.1 & 0.0 $\pm$ 0.0 & 0.0 $\pm$ 0.0 & 0.12 \\
ALE & 82.3 $\pm$ 0.5 & 35.4 $\pm$ 1.2 & 49.5 $\pm$ 0.8 & 0.51 \\
SJE & 81.1 $\pm$ 0.6 & 33.1 $\pm$ 1.4 & 47.0 $\pm$ 1.1 & 0.48 \\
DeVise & 83.5 $\pm$ 0.4 & 38.2 $\pm$ 0.9 & 52.4 $\pm$ 0.7 & 0.54 \\
Attribute-GAN & 86.4 $\pm$ 0.8 & 58.2 $\pm$ 1.5 & 69.5 $\pm$ 1.2 & 0.71 \\
WGAN-ZSL & 87.1 $\pm$ 0.7 & 61.4 $\pm$ 1.3 & 72.0 $\pm$ 1.0 & 0.73 \\
LisGAN & 89.4 $\pm$ 0.6 & 67.5 $\pm$ 1.1 & 76.9 $\pm$ 0.9 & 0.78 \\
CADA-VAE & 91.2 $\pm$ 0.4 & 75.2 $\pm$ 0.8 & 82.4 $\pm$ 0.6 & 0.83 \\
\textbf{Diff-LM-GZSL} & \textbf{91.0 $\pm$ 0.3} & \textbf{87.5 $\pm$ 0.5} & \textbf{89.2 $\pm$ 0.4} & \textbf{0.90} \\ \hline
\end{tabular}
\end{table}

The cross-load generalization ability of Diff-LM-GZSL is particularly noteworthy. While conventional methods suffer a sharp decline in unseen accuracy when the load conditions shift, our model maintains consistent performance. This is attributed to the language-driven semantic module, which provides load-invariant descriptors such as ``localized pitting on inner race'' and ``periodic impulsive components,'' allowing the diffusion model to focus on the intrinsic physical signatures of the faults rather than load-specific noise.

To further validate the contribution of each module, a comprehensive ablation study was conducted on the CWRU dataset \cite{Downs1993TEP}, as summarized in Table \ref{tab:cwru_ablation}. The removal of the LLM-based linguistic module (leaving only binary attributes) leads to a 12.5% drop in $H$, highlighting the superiority of rich textual descriptions over simplistic binary vectors. Replacing the latent diffusion model with a standard WGAN results in a 9.8% decrease in accuracy, confirming the diffusion model's superior capability in capturing complex multi-modal distributions. Furthermore, the exclusion of the Max-Mean Discrepancy (MMD) alignment loss causes the model to struggle with domain shift, as the generated features fail to overlap perfectly with real feature distributions in the latent space.

\begin{table}[ht]
\centering
\caption{Ablation Study on CWRU Dataset (Case III)}
\label{tab:cwru_ablation}
\begin{tabular}{lccc}
\hline
Configuration & $acc\_s$ (\%) & $acc\_u$ (\%) & $H$ (\%) \\ \hline
Full Diff-LM-GZSL & 91.0 & 87.5 & 89.2 \\
w/o LLM Semantic (Binary only) & 86.4 & 68.9 & 76.7 \\
w/o Diffusion (WGAN instead) & 88.5 & 72.1 & 79.4 \\
w/o MMD Alignment & 89.2 & 78.4 & 83.5 \\
\hline
\end{tabular}
\end{table}

Hyperparameter sensitivity analysis was performed to ensure optimal performance. Fig. 7 presents the confusion matrix for the CWRU dataset, showing that most misclassifications occur between faults of the same type but different diameters, which is expected given the similarity in their spectral signatures. Fig. 9 illustrates the loss convergence curves, indicating that the latent diffusion training stabilizes after approximately 800 epochs. The influence of the diffusion steps $T$ was also investigated. We observed that increasing $T$ from 500 to 1500 progressively improves the generation quality; however, $T=2000$ yields marginal gains while significantly increasing the computational overhead. Thus, $T=1000$ was selected as the optimal trade-off. The alignment loss weight $\alpha$ and the number of generated samples $N$ were found to be critical for the diversity of the synthesized features, with $N=400$ providing sufficient coverage of the unseen class manifolds. Finally, Fig. 11 compares the quality of features generated by Real sampling, GAN-based methods, and Diff-LM-GZSL, where the t-SNE visualization confirms that our method produces clusters that are more cohesive and better separated.

\subsection{Discussion and Time Complexity}
The experimental results across the TEP, Hydraulic, and CWRU datasets reveal consistent trends that underscore the efficacy of the proposed Diff-LM-GZSL framework. A primary finding is that the language-driven semantic alignment consistently outperforms traditional binary attribute-based ZSL. By utilizing an LLM to extract granular expert knowledge as semantic embeddings, we bridge the gap between low-level signal features and high-level diagnostic concepts more effectively than rigid attribute labels. This is particularly evident in industrial scenarios where fault symptoms are nuanced and difficult to quantify with binary tags. 

Furthermore, the latent diffusion model demonstrates a remarkable ability to generate high-fidelity synthetic features for unseen classes. Unlike GANs, which often suffer from mode collapse and training instability, the diffusion process provides a more robust likelihood-based estimation of the data manifold. This results in synthetic samples that better represent the intra-class variance of faults, leading to a significant reduction in the bias towards seen classes during GZSL classification.

Computational complexity is an essential consideration for industrial deployment. Table \ref{tab:time_complexity} compares the training and inference times of Diff-LM-GZSL against representative GAN and VAE based methods. While the training phase of the diffusion model is more time-consuming (approx. 4.2 hours on an NVIDIA RTX 3090) due to the iterative denoising process, the inference time remains highly competitive. The generation of synthetic features for classification is a one-time process during the training/adaptation phase; once the classifier is trained, the real-time diagnosis of incoming signals is nearly instantaneous ($\sim$12 ms per sample), making it suitable for online monitoring.

\begin{table}[ht]
\centering
\caption{Computational Complexity Comparison}
\label{tab:time_complexity}
\begin{tabular}{lcc}
\hline
Method & Training Time (h) & Inference Latency (ms) \\ \hline
WGAN-ZSL & 1.5 & 10 \\
LisGAN & 2.1 & 11 \\
CADA-VAE & 0.8 & 8 \\
\textbf{Diff-LM-GZSL} & 4.2 & 12 \\ \hline
\end{tabular}
\end{table}

Despite its strengths, certain limitations persist. The method's performance is intrinsically linked to the quality of the textual descriptions provided to the LLM; ambiguous or overly technical descriptions that deviate from the LLM’s pre-training data may lead to sub-optimal embeddings. Additionally, although the latent space approach mitigates some of the diffusion model's slowness, the iterative nature of $T$ steps still poses a bottleneck for applications requiring ultra-fast adaptation to new fault types.

\subsection{Extended Analysis: Impact of Signal Processing Techniques}
To assess the sensitivity of Diff-LM-GZSL to the initial transformation of 1D vibration signals, we expanded the evaluation on the CWRU dataset by comparing STFT, CWT, and Wigner-Ville Distribution (WVD). While STFT is computationally efficient, its fixed window size limits its resolution. CWT, used in our main results, provides a more adaptive time-frequency tiling, allowing the diffusion model to better capture short-duration impulsive fault signatures. Surprisingly, WVD, despite providing high resolution, introduced cross-term interference that negatively impacted the diffusion model's ability to learn clean latent representations, leading to a 3.4\% drop in the unseen class $acc\_u$. This suggests that for latent diffusion-based GZSL, the clarity and physical interpretability of the 2D input are as critical as the resolution itself.

Furthermore, we investigated the resilience of the linguistic semantic module under different signal-to-noise ratios (SNRs). Industrial environments are inherently noisy, and the ability of a ZSL model to generalize from noise-free descriptions to noisy real-world data is paramount. We simulated additive Gaussian noise at SNR levels of 0 dB, -5 dB, and -10 dB. Diff-LM-GZSL maintained an $H$ value above 75\% even at -5 dB SNR, whereas GAN-based methods plummeted below 60\%. This robustness stems from the LLM's high-level semantic descriptions, which act as a robust prior, effectively guiding the diffusion process to reconstruct the "clean" fault manifold even when the latent space is perturbed by noise.

\subsection{Statistical Significance and Cross-Validation}
To ensure that the performance gains of Diff-LM-GZSL over benchmarks like LisGAN and CADA-VAE are statistically significant, we performed a 5-fold cross-validation on the CWRU dataset. In each fold, a different subset of three fault types was held out as unseen. The results, shown in Table \ref{tab:cross_val}, indicate that our method consistently maintains its lead. A paired t-test between Diff-LM-GZSL and the second-best method, CADA-VAE, resulted in a p-value of 0.003, which is well below the 0.05 threshold for statistical significance. This confirms that the observed improvements are not due to a specific choice of seen/unseen splits but are inherent to the model's architecture and training objective.

\begin{table}[ht]
\centering
\caption{5-Fold Cross-Validation on CWRU Dataset (Unseen Split)}
\label{tab:cross_val}
\begin{tabular}{lccccc}
\hline
Method & Fold 1 & Fold 2 & Fold 3 & Fold 4 & Fold 5 \\ \hline
CADA-VAE & 82.4 & 81.1 & 83.5 & 80.9 & 82.2 \\
\textbf{Diff-LM-GZSL} & \textbf{89.2} & \textbf{88.7} & \textbf{90.1} & \textbf{88.4} & \textbf{89.5} \\ \hline
\end{tabular}
\end{table}

The variance across folds was only 0.65\%, compared to 1.12\% for the GAN-based baselines. This stability is a direct consequence of the diffusion model's likelihood-based training, which avoids the equilibrium-finding instabilities associated with adversarial training. By providing a stable generative foundation, Diff-LM-GZSL offers a more reliable tool for industrial health monitoring, where the cost of misclassification can be significant.

\subsection{Interpretability of Latent Diffusion Features}
One of the most compelling aspects of using language-driven semantics is the interpretability it brings to the generated features. By examining the cross-attention maps within the latent diffusion module, we can visualize which words in the fault description (e.g., "outer race", "periodic", "600Hz") activate specific regions in the generated feature vectors. Our analysis shows that words related to the physical location of the fault ("inner" vs "outer") correspond to low-frequency components in the latent space, while words describing the severity ("pitting", "spalling") correlate with high-frequency variations. This mapping confirms that the model is not merely memorizing patterns but is learning a physically grounded relationship between linguistic expertise and signal characteristics, fulfilling the "semantic alignment" goal of our research.
